\documentclass[11pt]{article}
\usepackage{icmlko}
\begin{document}

\title{의존도를 계기로 만들었더니 행 수가 아니라고 한다 --- 흔들리려면
흔들 신호가 있어야 한다}
\author{월드모델 실험실}
\date{노트 453 · 2026년 8월 2일}
\maketitle

\begin{abstract}
노트 452 의 팝업 $-0.3089$ 는 \emph{우연히 나온} 의존도 측정이었다 ---
자기 학습 행이 한 줄도 안 바뀌었는데 점수가 움직였다. 이번엔 대놓고
계기로 만들었다. 작은 도메인 넷의 학습 행을 \textbf{고정}하고 빌려주는
일곱만 되뽑기 열 번, 고정된 넷의 점수 SD 를 읽는다 --- 그 흔들림은
\textbf{순수하게 남 때문}이다.

등록한 문장은 통과했다. \textbf{팝업 $0.0349$ 대 시장팝업 $0.0067$ ---
$5.2$배.}

\textbf{그런데 그 뒤에 있던 이야기는 죽었다.}

\begin{center}
\begin{tabular}{lrrr}
\hline
고정 도메인 & 학습 행 & 의존도 SD & 같은 시기 천장 \\
\hline
팝업 & $24$ & $0.0349$ & $+0.2884$ \\
도서 & $80$ & $0.0398$ & $+0.4724$ \\
\textbf{아이돌} & $122$ & $\mathbf{0.0447}$ & $+0.5932$ \\
\textbf{시장팝업} & $123$ & $\mathbf{0.0067}$ & $\mathbf{-0.0901}$ \\
\hline
\end{tabular}
\end{center}

\textbf{아이돌과 시장팝업은 학습 행이 한 줄 차이인데 의존도가 $6.7$배
다르다.} 행 수와 SD 의 스피어만은 $-0.200$ --- \emph{``적게 배운
도메인일수록 남에게 휘둘린다''는 그림은 틀렸다.}
\end{abstract}

\begin{figure}[t]
\centering
\includegraphics[width=\textwidth]{figs/f1.pdf}
\caption{\textbf{가} --- 도메인별 의존도. 붉은 넷만 자기 행이 안 흔들려서
순수한 값이다(회색 일곱은 자기 행도 흔들려 섞였다). \textbf{나} ---
고정된 넷을 노트 445 의 천장 위에 놓았다. 천장이 $0$ 아래인 시장팝업만
안 흔들린다.}
\end{figure}

\section{그러면 무엇이 정하나}

넷뿐이라 맞추기는 못 하지만, 갈라지는 자리는 분명하다. \textbf{시장팝업만
같은 시기 천장이 음수}($-0.0901$, 노트 445)다. 자기 축으로 자기 라벨을
교차검증하면 $0$ 보다 못하다.

\textbf{흔들리려면 흔들 신호가 있어야 한다.} 시장팝업의 점수는
$+0.1526$ 에 붙박여 있고 그 바닥은 자기 자료가 아니라 판 전체의 평균적인
경향에서 온다. 빌려주는 쪽을 되뽑기로 흔들어도 \emph{그 평균은 잘 안
움직이므로} 시장팝업도 안 움직인다. 반대로 아이돌은 천장 $+0.5932$ 에
실제 점수 $+0.4752$ --- 움직일 여지가 크고, 그 여지의 상당 부분을 남에게
기대고 있다.

노트 445 는 시장팝업을 두고 \emph{``가진 게 전부 빌린 것''}이라고 적었다.
이 노트가 그 문장을 고친다 --- \textbf{빌린 것이 전부인 것과, 빌린 것이
많은 것은 다르다.} 시장팝업은 \emph{적게} 빌렸는데 그게 전부인 것이고,
아이돌은 \emph{많이} 빌렸다.

\section{계기는 남는다}

이야기는 죽었지만 \textbf{계기는 쓸 만하다}. 자기 행을 고정하고 남을
흔드는 것은 값이 싸고(되뽑기 열 번 = 열 번 적합) 뜻이 분명하다. 노트 452
가 만든 조항 --- \emph{학습 구성을 바꾸는 변경은 도메인별로 자기 행이
얼마나 바뀌었는지 같이 적는다} --- 과 짝이 된다. 이제 그 조항을 쓸 때
\textbf{$0\%$ 바뀐 도메인의 변화가 큰지 작은지를 미리 안다.}

\section*{남은 것}

\begin{itemize}
\item \textbf{일곱 도메인의 순수 의존도.} 지금은 자기 행도 흔들려 섞여
있다. 하나씩 고정하면 열한 번 더 돌려야 한다 --- 값이 있나.
\item \textbf{의존도와 노트 445 의 모음 값}(판 $-$ 혼자)이 같은 것을
재나. 넷으로는 못 보지만 열하나면 볼 수 있다.
\item \textbf{시장팝업의 천장이 왜 음수인가} --- 세 노트째 같은 자리를
가리킨다(노트 444 의 CN 과 같은 태거 갈래).
\end{itemize}

\end{document}
